\Subhead{1.7.2}{Деление частоты}
Вырожденный параметрический генератор в устойчивом синхронном режиме колеблется на $\omega_p/2$. Если собственная частота немного отличается, меняется относительная фаза и амплитуда, а частота остаётся половиной накачки в пределах области захвата. Поэтому устройство служит делителем частоты на два.

Уравнение инвариантно при $x\to-x$: два устойчивых фазовых состояния отличаются на $\pi$, а для накачки этот сдвиг равен целому периоду. Слабая затравка может выбрать фазу. Это принцип параметрона. Относительная частотная стабильность наследуется от накачки при сохранении синхронизма; фазовый шум прибора и срывы деления остаются возможны. Более высокие целочисленные деления требуют других нелинейных резонансов или каскада; из двухмодовых линейных уравнений они не следуют.
